% 洛伦兹群（综述）
% license CCBYSA3
% type Wiki

本文根据 CC-BY-SA 协议转载翻译自维基百科\href{https://en.wikipedia.org/wiki/Lorentz_group}{相关文章}

在物理学和数学中，洛伦兹群是指闵可夫斯基时空上的所有洛伦兹变换所构成的群。闵可夫斯基时空是所有（非引力）物理现象的经典与量子背景。洛伦兹群的名字来源于荷兰物理学家亨德里克·洛伦兹。

例如，以下定律、方程和理论都遵循洛伦兹对称性：
\begin{itemize}
\item 狭义相对论的运动学定律
\item 电磁学理论中的麦克斯韦场方程
\item 电子理论中的狄拉克方程
\item 粒子物理学的标准模型
\end{itemize}
洛伦兹群体现了自然界所有已知基本规律中时空的基本对称性。在时空区域足够小、引力差异可以忽略的情况下，物理规律与狭义相对论一样，保持洛伦兹不变性。
\subsection{基本性质}
洛伦兹群是庞加莱群（即闵可夫斯基时空所有等距变换所成的群）的一个子群。洛伦兹变换严格来说就是那些保持原点不变的等距变换。因此，洛伦兹群是闵可夫斯基时空等距群关于原点的各向同性子群。出于这个原因，洛伦兹群有时被称为齐次洛伦兹群，而庞加莱群则被称为非齐次洛伦兹群。洛伦兹变换是线性变换的一个例子；而闵可夫斯基时空的一般等距变换则是仿射变换。
\subsubsection{物理学定义}
\begin{figure}[ht]
\centering
\includegraphics[width=6cm]{./figures/c8d228c33aac897c.png}
\caption{亨德里克·安东·洛伦兹（Hendrik Antoon Lorentz, 1853–1928），洛伦兹群即以他的名字命名。} \label{fig_LuoLz_1}
\end{figure}
假设有两个惯性参考系 $(t, x, y, z)$ 与 $(t', x', y', z')$，以及两个点 $P_1, P_2$。洛伦兹群是所有在这两个参考系之间的变换集合，这些变换保持光在两点之间传播的速度不变：
\[
c^{2}(\Delta t')^{2}-(\Delta x')^{2}-(\Delta y')^{2}-(\Delta z')^{2}
=
c^{2}(\Delta t)^{2}-(\Delta x)^{2}-(\Delta y)^{2}-(\Delta z)^{2}.~
\]
用矩阵形式表示，这些变换是所有满足下式的线性变换 $\Lambda$：
\[
\Lambda^{\mathsf{T}} \eta \Lambda = \eta,
\qquad
\eta = \operatorname{diag}(1,-1,-1,-1).~
\]
\subsubsection{数学定义}
在数学上，洛伦兹群可以描述为不定正交群$O(1,3)$，即保持如下二次型的不变的矩阵李群：
\[
(t,x,y,z)\;\mapsto\; t^{2}-x^{2}-y^{2}-z^{2}~
\]
作用在$\mathbb{R}^4$上（赋予这一二次型的向量空间有时记作$\mathbb{R}^{1,3}$）。当写成矩阵形式时（参见经典正交群），这一二次型在物理学中被解释为闵可夫斯基时空的度规张量。
\subsubsection{符号说明}
在表示洛伦兹群时，$O(1,3)$与$O(3,1)$两种记号都被广泛使用。前者指保持一个符号型为$(+,-,-,-)$ 的度规的矩阵，后者则对应符号型为 $(-,+,+,+)$ 的度规。由于度规整体符号在定义方程中无关紧要，因此这两类矩阵群是相同的。

在当代，有部分领域倾向采用$(1,3)$的记号，而$(3,1)$的记号在现今实践和历史文献中依然大量存在。本文所描述的一切内容，对$O(3,1)$的记号同样成立，只需作相应的符号替换。类似的考虑同样适用于相关定义（例如$SO^{+}(1,3)$与$SO^{+}(3,1)$）。
\subsubsection{数学性质}
洛伦兹群是一个六维的、非紧致的、非阿贝尔的实李群，并且它不是连通的。其四个连通分支都不是单连通的。[1] 洛伦兹群的单位元所在的连通分支本身构成一个群，通常称为受限洛伦兹群，记作$SO^{+}(1,3)$。受限洛伦兹群由那些同时保持空间取向与时间方向的洛伦兹变换组成。其基本群的阶为 2，而它的普适覆盖群——不定自旋群$Spin(1,3)$——同构于特殊线性群$SL(2,\mathbf{C})$和辛群$Sp(2,\mathbf{C})$。这些同构使得洛伦兹群可以作用于大量对物理学非常重要的数学结构，最典型的就是旋量。因此，在相对论性量子力学和量子场论中，通常把$SL(2,\mathbf{C})$称作洛伦兹群，并理解为$SO^{+}(1,3)$是它的一个特定表示（即矢量表示）。

洛伦兹群在闵可夫斯基空间上的一种常见表示使用了双四元数，它们构成一个合成代数。洛伦兹变换的等距性来源于这种合成性质：
\[
|pq| = |p| \times |q|.~
\]
洛伦兹群的另一性质是共形性，即角度保持。洛伦兹推动作用为时空平面上的双曲旋转，而这种“旋转”保持双曲角，即相对论中使用的快度的度量。因此，洛伦兹群是时空共形群的一个子群。

需要注意的是，本文称$O(1,3)$为“洛伦兹群”，$SO(1,3)$为“正规洛伦兹群”，而 $SO^{+}(1,3)$为“受限洛伦兹群”。许多作者（尤其在物理学中）使用 “Lorentz group” 一词指代$SO(1,3)$（甚至有时是$SO^{+}(1,3)$）而不是 $O(1,3)$。在阅读此类文献时，需要明确辨清他们所指的具体是哪一个群。
\subsubsection{连通分支}
\begin{figure}[ht]
\centering
\includegraphics[width=6cm]{./figures/239174c7897a6114.png}
\caption{二维空间加一个时间维度中的光锥} \label{fig_LuoLz_2}
\end{figure}
因为洛伦兹群$O(1,3)$是一个李群，所以它既是一个群，同时也可以在拓扑学上描述为一个光滑流形。作为流形，它有四个连通分支。直观上，这意味着它由四个在拓扑上彼此分离的部分组成。

这四个连通分支可以通过其元素的两类变换性质来分类：
\begin{itemize}
  \item 一些元素在时间反演的洛伦兹变换下会被反转，例如，一个未来指向的类时向量会被反转为过去指向的向量；
  \item 一些元素在非正规洛伦兹变换下会导致取向被反转，例如某些vierbein（四重基，或称四重标架）。
\end{itemize}
保持时间方向的洛伦兹变换称为\textbf{正时}。正时变换的子群通常记作$O^{+}(1,3)$。保持空间取向的称为\textbf{正规}变换，作为线性变换，它们的行列式为 $+1$。（而非正规洛伦兹变换的行列式为$-1$。）正规洛伦兹变换的子群记作$SO(1,3)$。[a]

保持空间取向与时间方向的所有洛伦兹变换所构成的子群称为\textbf{正规正时洛伦兹群}，或称为\textbf{受限洛伦兹群}，记作$SO^{+}(1,3)$。

这四个连通分支的集合可以通过商群$O(1,3)\,/\,SO^{+}(1,3)$赋予群结构，该商群同构于克莱因四元群。$O(1,3)$中的任意元素都可以表示为一个正规正时变换与离散群元素的半直积：
\[
\{1,\, P,\, T,\, PT\},~
\]
其中$P$和$T$分别是宇称与时间反演算子：
\[
P = \operatorname{diag}(1,-1,-1,-1), 
\quad 
T = \operatorname{diag}(-1,1,1,1).~
\]
因此，一个任意的洛伦兹变换可以由一个正规正时洛伦兹变换，加上两个额外的比特信息来唯一确定，这两个比特信息决定了四个连通分支中的哪一个。这一模式是有限维李群的典型特征。
\subsection{受限洛伦兹群}
受限洛伦兹群$SO^{+}(1,3)$是洛伦兹群的单位元分支，这意味着它包含了所有可以通过一条连续曲线与单位元相连的洛伦兹变换。受限洛伦兹群是完整洛伦兹群的一个连通正规子群，并且具有相同的维数，在此情形下为六维。

受限洛伦兹群由普通的空间旋转和洛伦兹推动生成（洛伦兹推动是包含类时方向的双曲空间中的“旋转”[2]）。由于每一个正规正时洛伦兹变换都可以写作一次旋转（由三个实参数指定）与一次推动（也由三个实参数指定）的乘积，因此需要六个实参数来唯一确定一个任意的正规正时洛伦兹变换。这也是理解受限洛伦兹群为何是六维的一种方式。（另见洛伦兹群的李代数。）

所有旋转构成的集合形成一个李子群，与普通旋转群$SO(3)$同构。而所有推动构成的集合并不构成子群，因为两个推动相乘一般并不再是一个推动。（相反，两个非共线推动等价于一次推动加一次旋转，这与托马斯旋转相关。）沿某个方向的推动，或绕某个轴的旋转，生成一个一参数子群。
\subsubsection{传递性曲面}
\begin{figure}[ht]
\centering
\includegraphics[width=10cm]{./figures/3e8fbc86e2952080.png}
\caption{} \label{fig_LuoLz_3}
\end{figure}
如果一个群$G$作用在空间$V$上，那么一个曲面$S \subset V$称为\textbf{传递性曲面}，如果$S$在$G$下是不变的（即$\forall g \in G,\, \forall s \in S:\, gs \in S$），并且对于任意两点$s_{1}, s_{2} \in S$，存在某个$g \in G$使得$g s_{1} = s_{2}$。  

根据洛伦兹群的定义，它保持如下二次型：
\[
Q(x) = x_{0}^{2} - x_{1}^{2} - x_{2}^{2} - x_{3}^{2}.~
\]
在平直时空 $\mathbb{R}^{1,3}$ 上，正时洛伦兹群 $O^{+}(1,3)$ 的传递性曲面 $Q(x)=\text{const.}$ 如下所示：[3]
\begin{itemize}
  \item $Q(x) > 0,\, x_{0} > 0$：这是双叶双曲面（hyperboloid of two sheets）的上支。其上的点与原点相隔一个未来类时向量。
  \item $Q(x) > 0,\, x_{0} < 0$：这是该双曲面的下支。其上的点是过去类时向量。
  \item $Q(x) = 0,\, x_{0} > 0$：这是光锥的上支，即未来光锥。
  \item $Q(x) = 0,\, x_{0} < 0$：这是光锥的下支，即过去光锥。
  \item $Q(x) < 0$：这是单叶双曲面（hyperboloid of one sheet）。其上的点与原点是类空间分离的。
  \item 原点：$x_{0} = x_{1} = x_{2} = x_{3} = 0$。
\end{itemize}
这些曲面都是三维的，因此相应的图像并不完全忠实，但对于$O^{+}(1,2)$的对应情形则是忠实的。对于完整的洛伦兹群而言，传递性曲面只有四类，因为变换$T$会将双曲面（或光锥）的上支映射到下支，反之亦然。

\textbf{对称空间的视角}

对上述传递性曲面的另一种等价表述方式是将其看作李理论意义下的\textbf{对称空间}。例如，根据轨道–稳定子定理，双叶双曲面的上支可以写作商空间SO^{+}(1,3) \,/\, SO(3).
\]
此外，这一上支还可以作为三维双曲空间的一个模型。
